\appendix

\section{LLM Judge System Prompt}
\label{app:judge_prompt}

The following system prompt was used to instruct the GPT-4o judge during all evaluation runs. The prompt explicitly constrains the model to a factual verification task, rejecting responses that apologize, hallucinate, or miss critical details.

\begin{quote}
\small
\textit{You are an expert evaluation judge. Your task is to determine whether an AI assistant successfully answered a probe question by recalling a specific fact from earlier in the conversation.}

\textit{Return preserved=true if the assistant's answer correctly incorporates the Expected Fact (even if paraphrased slightly or formatted differently). Return preserved=false if the assistant apologizes, says it doesn't know, hallucinates a different fact, or misses critical details.}
\end{quote}

\section{Extended Results: Per-Run CPR Stability}
\label{app:extended_results}

Table~\ref{tab:app_cpr_runs} presents the unaggregated Continuity Preservation Rate (CPR) for each of the five independent runs comprising the $N=750$ evaluation suite. Each run evaluates $n=150$ failover events per system. The treatment CPR demonstrates tight stability, varying by a maximum of 0.6 percentage points across all runs.

\begin{table}[h]
\centering
\caption{Continuity Preservation Rate broken down by individual run. Confidence intervals are 95\% Wilson score intervals.}
\label{tab:app_cpr_runs}
\small
\begin{tabular}{ccrccrc}
\toprule
 & \multicolumn{3}{c}{\textbf{Baseline (Stateless)}} & \multicolumn{3}{c}{\textbf{Treatment (History-Forwarding)}} \\
\cmidrule(lr){2-4} \cmidrule(lr){5-7}
\textbf{Run} & \textbf{Preserved} & \textbf{CPR (\%)} & \textbf{95\% CI} & \textbf{Preserved} & \textbf{CPR (\%)} & \textbf{95\% CI} \\
\midrule
1 & 0 / 150 & 0.0\% & [0.0, 2.5] & 149 / 150 & 99.3\% & [96.3, 99.9] \\
2 & 0 / 150 & 0.0\% & [0.0, 2.5] & 149 / 150 & 99.3\% & [96.3, 99.9] \\
3 & 0 / 150 & 0.0\% & [0.0, 2.5] & 149 / 150 & 99.3\% & [96.3, 99.9] \\
4 & 0 / 150 & 0.0\% & [0.0, 2.5] & 148 / 150 & 98.7\% & [95.3, 99.6] \\
5 & 0 / 150 & 0.0\% & [0.0, 2.5] & 149 / 150 & 99.3\% & [96.3, 99.9] \\
\midrule
\textbf{Pooled} & \textbf{0 / 750} & \textbf{0.0\%} & \textbf{[0.0, 0.5]} & \textbf{744 / 750} & \textbf{99.2\%} & \textbf{[98.3, 99.6]} \\
\bottomrule
\end{tabular}
\end{table}

\section{Sample Failover Transcripts}
\label{app:transcripts}

The following transcripts are sampled directly from the manual audit of Run 5. They illustrate the History-Forwarding mechanism in action: the fallback provider (Anthropic) correctly answers a highly specific probe question using context originally established with the primary provider (OpenAI). 

All conversations are synthetic and contain no personally identifiable information.

\paragraph{Case 1: Invented Term (conv-093)}
\begin{quote}
\small
\textbf{Turn 0 (User):} In the game, the currency is called Plexarine. \\
\textit{[... intervening filler dialogue processed by OpenAI ...]} \\
\textbf{Turn 10 (User, Probe):} What is the currency called in the game I described? \\
\textcolor{red}{\ding{55}} \textit{[Primary Provider (OpenAI) fails at Turn 10]} \\
\textcolor{green!60!black}{\ding{51}} \textit{[Proxy forwards full history to Fallback Provider (Anthropic)]} \\
\textbf{Turn 10 (Fallback Asst):} The currency in the game you described is called \textbf{Plexarine}.
\end{quote}

\paragraph{Case 2: Specific Preference (conv-027)}
\begin{quote}
\small
\textbf{Turn 0 (User):} I prefer Gitflow with squash merges for version control workflows. \\
\textit{[... intervening filler dialogue processed by OpenAI ...]} \\
\textbf{Turn 8 (User, Probe):} What version control workflow did I say I prefer? \\
\textcolor{red}{\ding{55}} \textit{[Primary Provider (OpenAI) fails at Turn 8]} \\
\textcolor{green!60!black}{\ding{51}} \textit{[Proxy forwards full history to Fallback Provider (Anthropic)]} \\
\textbf{Turn 8 (Fallback Asst):} You said you prefer Gitflow with squash merges.
\end{quote}

\paragraph{Case 3: Numeric Value (conv-080)}
\begin{quote}
\small
\textbf{Turn 0 (User):} The new protocol we're designing is codenamed Syndexium. \\
\textit{[... intervening filler dialogue processed by OpenAI ...]} \\
\textbf{Turn 6 (User, Probe):} What codename did I give the new protocol? \\
\textcolor{red}{\ding{55}} \textit{[Primary Provider (OpenAI) fails at Turn 6]} \\
\textcolor{green!60!black}{\ding{51}} \textit{[Proxy forwards full history to Fallback Provider (Anthropic)]} \\
\textbf{Turn 6 (Fallback Asst):} You called it Syndexium.
\end{quote}

\section{Harness Configuration Details}
\label{app:config}

To facilitate full reproduction, the exact configuration parameters used for the $N=750$ evaluation run are provided below.

\begin{itemize}
    \item \textbf{Concurrency ($C$):} 100 simultaneous simulated sessions, managed via \texttt{asyncio.Semaphore}.
    \item \textbf{Primary Provider:} OpenAI (\texttt{gpt-4o-mini}), \texttt{temperature=0.0}.
    \item \textbf{Fallback Provider:} Anthropic (\texttt{claude-3-5-sonnet-20240620}), \texttt{temperature=0.0}.
    \item \textbf{Retry Strategy:} Exponential backoff with jitter on HTTP 429, 502, 503, 504.
    \item \textbf{Backoff Formula:} $t_{\text{wait}} = \min(30\text{s}, 2^a + \mathcal{U}(0, 1))$ where $a$ is the attempt index.
    \item \textbf{Maximum Retries:} 5 attempts per failover request.
    \item \textbf{Proxy Read Timeout:} 300.0 seconds (to accommodate the Anthropic SDK's native backoff behavior without dropping the client connection).
    \item \textbf{Fault Injection Point:} ``Late'' mode (failure injected exclusively on the final probe turn).
\end{itemize}
